\section{形态计算之后：身体封装与演化摊销}

\subsection{物理智能的残差问题与身体封装的三层结构}

承受总要落在某个具体结构上。生命把大量残差修正压进了演化、发育和环境的共适配，机器没有这套结构，它在物理世界中的每一次行动，都要靠显式工程结构来吸收偏差。第三章留下的压缩判据决定了模型保留哪些变量，但这些变量在物理接触中由谁修正、以多大延迟修正、失败之后能否恢复，并不是模型内部可以回答的问题。

物理智能最容易被误解的一点，是把身体看成智能的外设。在这种理解里，真正的智能位于大脑、模型或控制系统之中，身体只是执行命令的机械端口，系统先理解世界，再规划动作，最后由身体把计划落实到现实；按照这种图式，只要模型足够强、传感器足够多、控制足够精细，机器就应当能够像人一样进入现实世界。但物理智能的真实困难，不在于生成一个动作，而在于在开放现实中连续行动，而连续行动要先承认一件模型不愿承认的事：现实会持续产生模型没有预写的偏差。

\subhead{物理智能的残差问题}

现实不是等待执行的静态对象，而是一个不断产生残差的开放系统。模型可以给出动作，现实会给出偏差。地面可能比预期更滑，物体可能比视觉判断更重，杯壁可能有水，门把手可能松动，人的身体可能突然失衡，光线可能变化，传感器可能被遮挡，执行器可能存在延迟。每一个偏差单独看都很小，但物理任务不是一次问答，而是一条任务链：一次小的识别误差会改变抓取姿态，一次小的接触偏差会改变物体位置，一次定位漂移会进入下一轮规划，一次反馈延迟会放大控制误差，残差不会在动作之间自动清零，而会沿着任务链传播。

这里所说的\term{残差}，不是一个单一误差项，而是开放物理行动中未被当前模型、身体、控制器或运行体系完全吸收的偏差总称，它至少包括四类：模型残差，即世界模型或策略模型压缩后遗漏的状态变量；感知残差，即传感器噪声、遮挡、漂移和多模态不同步造成的状态估计偏差；接触残差，即摩擦、滑移、变形、冲击和力传递中的不可完全预测部分；运行残差，即异常处理、维护、验证、接管和责任追溯中的系统性负担。要紧的不是这些残差是否存在，它们在开放行动中总是存在，而是在连续任务链中它们是否会被及时吸收，还是沿接口传播并越过可恢复边界。因此，具身智能真正要解决的，不是能不能做某个动作，而是能不能在连续任务中把残差压在可恢复范围内。

抓杯子只是最小例子。现实中的机器人面对的是清洁房间、护理老人、整理桌面、搬运货物、工地协作、农田采摘、医院送药，每个任务都由许多动作串联而成，每个动作又改变下一个动作的初始条件，系统真正要证明的，不是它能不能成功一次，而是它能不能在长时间、多扰动、高责任的任务链中保持可恢复性。这里的困难有明确的门槛形状：超过某个残差阈值，能力不是平滑下降，而可能突然表现为停机、误抓、碰撞、失衡或请求人工接管。于是问题就落到承载残差的结构上，而人类身体的特殊之处，不是拥有完整世界模型，而是把大量残差修正压进了身体、经验和环境共适配之中。

\subhead{身体封装的三层结构}

拿起一个杯子时，人不会先计算杯子的精确几何、摩擦系数、重心位置、表面湿度、接触角度和力分布，人只是看一眼，伸手，触碰。手指立刻感到滑不滑，肌肉立刻感到重不重，关节立刻修正姿态，杯子一有微小滑移，指腹和肌肉就自动调整夹持力度。很多误差还没有进入意识，就已经被身体吸收，这就是\term{身体封装}。\key{身体不是智能的外设，而是残差管理系统。}它把开放现实中的大量偏差压缩在低延迟、连续反馈、自修正的闭环里。

身体封装不是说身体“很重要”这么简单，而是说身体本身是一套残差吸收结构。皮肤、肌肉、肌腱、关节、骨骼、前庭系统、本体感觉、痛觉和反射，并不是大脑之外的附属装置，而是感知、执行、缓冲、约束和修正的连续系统：手指柔软，会自动适配物体表面；肌肉有弹性，会吸收冲击；关节有限位，会限制危险动作；痛觉会标记过载；经验会提示哪些细节可以忽略，哪些细节意味着危险。身体不是把所有信息传给大脑再等待计算，而是在接触中直接处理一部分现实。

身体之所以能够封装残差，并不是因为它绕开了物理规律，而是因为它把一部分物理规律变成了行动结构。摆腿时的重力势能与动能转换、肌肉与肌腱中的弹性储能、指腹接触时的局部形变、关节限位对危险姿态的排除，都不是中央控制器事后补偿出来的结果，而是身体形态预先把某些误差限制在可承接的动力学通道中；成熟身体不是每一刻都用主动控制对抗重力、惯性、摩擦和接触，而是在许多场景中让这些物理关系先完成一部分约束、筛选和修正。由此看，形态计算的深层意义不只是“身体参与计算”，而是身体把一部分原本需要感知、估计和控制的残差，转化为材料、几何和被动动力学中的低成本约束\cite{paul2006}，被动动力行走机器人正是这条思路的直接证据：在恰当形态和动力学条件下，稳定运动并不必然依赖高强度中央控制，身体结构本身可以承担一部分运动组织\cite{mcgeer1990,full1999,collins2005}。

封装至少发生在三层。第一层是材料封装，皮肤、肌肉、肌腱和关节以物理方式吸收误差，柔软的指腹扩大接触面积，肌肉弹性缓冲冲击，步态中的身体弹性帮助稳定行走，这里的“计算”不是符号计算，而是材料和形态直接参与控制，形态计算研究之所以重要，正在于它把身体形状、材料顺应性、被动动力学和环境耦合纳入控制问题，一部分本来需要中央控制器处理的复杂性，可以由身体结构本身分担\cite{paul2006}。第二层是反馈封装，人通过视觉、触觉、本体感觉、前庭系统和肌肉张力，在行动中连续采样现实：地面是否湿滑，脚底一压就知道；物体是否松动，手指一碰就知道；对方身体是否失衡，手臂一接触就会调整，世界不需要被完整预装，而可以在行动中被局部询问，阻抗控制在机器人操作中成为基础概念，正是因为接触任务不能被还原为纯位置跟踪，操作者、执行器、物体和环境构成同一个动力学耦合系统，控制的对象不是孤立轨迹，而是交互关系\cite{hogan1985}。第三层是经验封装，人见过无数杯子、门把手、楼梯、湿地面、刀具、热锅、疲惫的人体和拥挤的空间，这些经验不是完整模型，而是大量可行动的有损压缩，人知道哪些细节通常可以忽略，哪些细节必须立即处理，身体让模型可以不完整，经验让行动可以不从零开始\cite{gibson1979}。

这一路径与 Moravec 悖论、具身认知和形态计算传统相通\cite{moravec1988,pfeifer2007,varela1991,gibsonpick2000,pfeiferbongard2006,hauser2011}。Moravec 悖论提示，对机器来说，抽象推理未必最难，反而是人类幼儿即可完成的感知-运动能力极难工程化；具身认知指出，智能不只存在于头脑中，而是在身体、环境和行动之间形成；形态计算进一步强调，身体材料与形态可以参与控制。本节要推进的是其中一层更具体的判断：身体封装的不是抽象计算量，而是残差的低延迟吸收能力。不过，这个判断很容易被读成身体天然优于机器，而原文自身也留下了修正条款：人类并不是只把复杂性内化，机器也不是只能把复杂性外化，人类把一部分残差交给环境，门把手、楼梯、桌椅、道路和工具，都是围绕人体尺度形成的外部封装，机器也可以把一部分残差压回身体，柔性夹爪、弹性关节、机械限位和顺应性材料，都是机器的内部封装。真正的差异不在于人内化、机器外化，而在于时间分母与基质可供性这两条更底层的轴，这一点留到本节最后一个小节处理；眼下要先回答的是另一个问题：既然机器也有办法修正残差，为什么它仍然显得沉重？

\subsection{机器复杂度展现与接口债务}

\subhead{机器复杂度展现与接口债务}

机器并不是没有办法修正残差，而且路径很多。它可以通过传感器展现残差，摄像头、深度相机、激光雷达、触觉、力矩、IMU、本体感觉，都可以让机器知道自己偏了多少、滑了没有、撞了没有、姿态漂了没有；可以通过控制算法修正残差，闭环控制、阻抗控制、模型预测控制、实时重规划，可以把小误差不断拉回轨道；可以通过身体结构吸收残差，柔性夹爪、欠驱动手、弹性关节、顺应性材料、机械限位、被动稳定结构，都可以让身体自己消化一部分误差\cite{mason1985,birglen2008}；可以通过更强模型预测残差，视觉-语言-动作模型、世界模型、模仿学习和策略模型，可以从较少输入中推断更多状态，提前避开危险动作\cite{shao2025,li2025}；可以通过环境改造减少残差，标准货架、二维码、专用夹具、固定路线、封闭园区、轨道系统、V2X、结构化仓库，都是让环境替机器降低开放性；还可以通过任务收缩和运行体系限制残差，扫地、仓储搬运、巡检、手术辅助、焊接、采摘、配送各自收缩任务边界，远程接管、安全员、异常工单、日志审计、维护人员和保险责任，则把残差转入更大的运行结构。真正的问题不在于办法够不够多，而在于人类把大量残差修正沉入身体、经验和环境，机器往往必须把残差修正展现为工程系统。

这就是\term{机器复杂度展现}。所谓复杂度展现，不是简单地“机器需要更多传感器”，也不是“机器硬件不如人”，它指的是：在人类身体中被连续、局部、低延迟处理的现实复杂性，在机器那里必须被拆成可测量、可同步、可控制、可验证、可维护、可追责的工程结构。人用手指感知滑移，机器可能需要触觉传感器、滑移检测算法、控制反馈和安全策略；人用肌肉缓冲碰撞，机器可能需要柔顺材料、力控策略和机械保护；人通过经验理解“轻一点”“别弄坏”“扶稳一点”，机器必须把这些模糊约束转化为动作规则、风险边界和责任记录。身体里原本连在一起的东西，被机器拆成了模块，模块一多，接口就出现，接口一多，残差传播路径就增加。

这可以称为\term{接口债务}：当隐含在身体闭环中的残差修正被拆成多个工程模块后，系统必须持续承担同步、标定、延迟补偿、故障传播、维护和验证的结构性负担。在人类身体中，感知、执行、缓冲、修正和安全边界不是五个独立系统，手指触到杯壁的瞬间，触觉、肌肉张力、关节角度、皮肤形变和经验判断已经在同一个身体闭环中完成局部修正，变量没有被完全命名，误差没有被完全展开，控制也没有被拆成彼此独立的模块；机器不同，滑移要被检测，触觉要被标定，视觉要和力觉同步，控制器要补偿延迟，执行器要承接误差，安全系统要判断边界，日志系统要记录行为，验证系统要证明可靠，责任系统要追溯后果。每拆出一个变量，就多一个接口，每多一个接口，就多一个同步问题、故障模式和维护点。若把模块间的耦合近似为两两都可能相关，接口数量的增长可以写成
\begin{equation}
N_{\text{if}}=\frac{n(n-1)}{2}
\end{equation}
其中 $n$ 是被显式拆分的模块数，$N_{\text{if}}$ 是系统需要维护的成对接口数，两者都是无量纲计数，系数 $1/2$ 来自成对通道无序，一条通道不重复计数。在每一对模块之间至少需要一条同步、标定或误差传递通道的假设下，这个式子给出的是接口数量的上限而非实计数：$n$ 从五增到十，上限从十增到四十五，增长快于模块数本身。真实系统的依赖通常不是全连接，很多成对通道并不活跃，实际接口数低于该上限，但方向不变，模块拆得越细，需要同步、标定和验证的关系增长得越快。展现因此不是把复杂性看清楚这么简单，而是把隐含复杂性转化为显式耦合结构。

接口债务也不是普通的集成成本。普通集成成本强调模块多了会麻烦，接口债务强调残差会沿接口传播：感知偏差会进入控制，控制偏差会进入执行，执行偏差会改变环境，环境变化又反过来污染下一次感知；模块越多，残差传播路径越多，任务链越长，残差越容易从局部误差变成系统偏移，机器的沉重不来自某一个部件，而来自这些显式结构之间不断累积的耦合债务。这里面还藏着一层更隐蔽的时间债务：机器系统中的视觉、触觉、力矩、本体感觉、规划器和执行器通常具有不同的采样率、延迟和相位，每个模块都在处理某个时间片上的世界，但这些时间片并不天然重合，控制器拿到的“当前状态”，往往已经是多路异步信号经过时间戳、滤波、外推和融合之后生成的工程对象。只要任务进入高速接触或连续操作，残差就不只是从一个模块传到另一个模块，还会从一个时间相位错位到另一个时间相位，系统修正的可能不是此刻的误差，而是刚刚过去的误差，下一次动作又会把这个滞后的修正写入新的物理状态；于是接口越多，系统不仅越难集成，也越难维持一个稳定、共同、可行动的“现在”。身体的接缝不需要被逐一命名，机器的接缝必须被逐一维护。

反过来看，“展现”也不等于“堆料”。成熟工程不一定是把所有东西做得更通用、更复杂，一个柔性夹爪可能比高精度机械手加复杂触觉阵列更适合抓软物体，一个固定工装可能比通用识别模型更可靠，一个机械限位可能比复杂软件安全策略更稳定，一个封闭园区可能比开放道路更容易自动驾驶，真正的问题不是传感器够不够多，而是残差由谁修正、在哪里修正、以什么物理和工程代价修正。软体机器人和欠驱动结构正是在这个方向上展开，它们不是简单模仿生物外形，而是试图让材料顺应性、接触几何和被动形变参与控制，使机器人不必把每一次接触都转化为高精度测量和中央计算；Rus 和 Tolley 的综述指出，柔性材料使机器人能够走出刚性机器人的设计、制造和控制路线，但软体结构也带来建模、感知、驱动和控制的新困难\cite{rus2015}。这说明身体可以吸收残差，但身体本身也会引入新的显式复杂性，它不是免费解法，而是另一种复杂性分配方式。既然复杂性总要被分配，把其中一部分交给更强的模型，就成了下一条看似最直接的出路。

\subsection{预测不能取消承受：任务链与断崖式失效}

\subhead{预测不能取消承受}

模型可以减少误差、避免危险接触，也能让现有硬件完成更多任务。物理闭环还需要足够的力、速度、感知与恢复能力：滑移发生以后，要有接触与控制结构接住；部件损坏以后，要有修复或替代路径。预测、执行和恢复的能力，应分别检查。

这并不是否认大脑的价值。更强的视觉模型、世界模型、动作模型和策略学习，确实可以让已有传感器和执行器做得更多，很多任务并不需要人类级身体，也不需要全套触觉皮肤和灵巧手，只要任务边界清楚、环境相对稳定、后果可控，一个好模型配合现有硬件就能形成可用闭环。但更强大脑只能推远能力边界，不能取消能力边界：它不能让执行器凭空获得力量，不能让材料凭空获得柔顺性，不能让传感器永不漂移，不能让开放世界的长尾风险无需验证。大脑越强，某些复杂性会从传感器和硬件转移到模型，但只要行动进入原子世界，残差仍然必须在某个位置被吸收，而在开放任务链中，残差通常只能被吸收、转移、隔离或延迟暴露，很难被简单消除。残差不会免费消失，只会换一个位置重新出现。

端到端模型、仿真、遥操作与新型传感都可能减少或消除具体故障，也可能把部分负担转移给数据、支持与维护。残差没有固定守恒量。评价时应当记录错误是否真正减少、出现位置怎样改变，以及长任务中的恢复成本。可检验的机制是残差何时传播并放大，而不是预先断言它只能转移。

\subhead{任务链与断崖式失效}

经验机制的第一条推论是，实验室里的短链条成功不能直接等同于开放世界中的长期可靠。一个机器人在标准桌面上抓取成功，不等于它能整理陌生家庭；仓储机器人能在标准货架间运行，不等于它能在医院走廊中处理人流、隐私、急救优先级和临时封路；自动驾驶能在规则道路上行驶，不等于它能稳定处理施工、逆光、雨雪、非标准手势和责任争议。现实任务会拉长链条、增加扰动、提高责任后果，并让微小误差在接口之间循环，因此残差机制并不意味着机器最终不能实现可靠具身智能，它反对的只是另一种理解：把开放世界具身智能看成可以无代价、无限平滑外推的能力曲线。

以“家庭整理桌面”为例，任务看似简单，实际却是一条高残差任务链。机器人首先要识别桌面物体，杯子、纸巾、药盒、钥匙、充电线、眼镜、半杯水、易碎摆件；视觉模型可能把透明杯壁、反光桌面和遮挡边界处理得足够好，但仍会留下状态残差，杯子是否装水、杯底是否潮湿、眼镜是否压在纸下、药盒是否空盒。下一步抓取会把状态残差转化为接触残差，夹爪角度略偏，杯身轻微旋转，纸张被带起，充电线缠住另一件物体；再下一步，接触残差会改变环境状态，杯子位置变了，纸堆塌了，原先可见的钥匙被遮挡，桌面布局不再等于上一轮规划的输入。此时系统面对的不是一个孤立错误，而是误差写入环境之后形成的新任务状态，若机器人缺少低延迟触觉、柔顺抓取、局部恢复和任务级回滚能力，一个小小的滑移就可能从“夹得不稳”升级为“物体掉落”，再升级为“液体洒出”“任务中断”“人工接管”。这也正是具身智能能力边界的准确位置：不是模型能否生成一个看似正确的动作，而是系统能否在一定任务长度、扰动强度和责任后果下，把残差压在可恢复范围内。

边界处的突然暴露，还与物理接触本身的多尺度结构有关。人的指腹、皮肤、肌肉、关节和本体感觉并不是在单一尺度上工作，微小纹理、局部滑移、接触面积变化、整体姿态调整和肌肉张力修正，常常在连续介质中同时发生，许多微小偏差在还没有变成显式错误之前，已经被低层材料和感觉-运动闭环吸收；机器系统则常常把这些尺度拆开，微观接触由材料承受，局部滑移由触觉或力传感器检测，整体姿态由视觉或状态估计算法修正，任务偏移再交给规划器处理。若某一尺度上的残差没有被及时吸收，它就可能跃迁到更高尺度：微小滑移变成抓取姿态改变，姿态改变变成物体状态改变，物体状态改变又污染下一轮感知与规划。所谓断崖式失效，很多时候并不是误差线性增大，而是残差跨过了某个尺度间隙，从局部扰动变成系统事件。断崖式失效是身体封装不足、尺度封装不足和接口债务过重的外在症状。

这个机制可以被证伪，而且证伪的方式很具体。如果未来大量开放世界机器人在家庭、医院、工地、农田和公共道路中，能够长期呈现平滑降级：错误可自我恢复，残差不沿任务链积累，人工接管率随扰动强度缓慢上升而不是突然跳升，维护和验证成本不随任务长度显著增加，那么本节关于“残差沿任务链传播并在边界处放大”的判断就应被降级，那将说明更强模型、身体结构、环境设计和运行体系已经把残差压回了可连续吸收的范围。现实中也已经出现了一些重要的技术信号，Bongard、Zykov 和 Lipson 展示了机器人通过连续自建模在损伤后恢复运动的可能性\cite{bongard2006}，Cully 等人进一步展示了机器人可以通过智能试错在很短时间内找到损伤后的补偿行为\cite{cully2015}；这些工作并不推翻残差机制，反而说明可靠具身智能必须拥有残差恢复结构。换言之，真正成熟的具身智能不是永不出错，而是能够在接近边界时优雅降级，它依赖的不是把误差消灭，而是把误差重新纳入可处理范围的能力。而这种能力在人类身上之所以显得毫不费力，原因不在当下这一刻，而在时间。

\subhead{演化摊销与残差管理}

人类行动的轻巧不是免费的，它的代价已经被极长时间尺度摊销。数亿年的生物演化塑造了皮肤、肌肉、骨骼、关节、神经反射和运动控制，无数个体的失败、损伤、死亡和繁衍，把不适合现实的身体结构淘汰掉；个体发育又通过婴儿期、童年和日常生活中的无数次摔倒、碰撞、抓握、试探和修正，把身体与现实重新校准；社会环境则围绕人的尺度设计了门把手、楼梯、桌椅、工具、厨房、道路、医院和工厂。人进入现实，并不是进入中性空间，而是进入一个与人类身体共同演化、共同设计、共同驯化过的生态位，这就是\term{演化摊销}。

演化摊销不是说人类身体没有代价，而是说这部分代价被分布在极长的时间尺度、海量个体试错、个体发育和社会环境改造之中，一个成人轻松拿起杯子，背后不是当下瞬间的简单计算，而是演化史、发育史、生活史和环境史共同沉淀出的身体封装，发育系统理论和生态心理学都强调，行动能力不是单纯从内部程序展开，而是在身体、环境和历史过程的耦合中形成\cite{thelen1994,clark1997}。到了这里，前面提到的那两条轴可以展开了。时间分母决定复杂性能否被长期摊销，人类身体和人类环境用演化、发育和文明设计摊销了极长时间，机器则试图在短工程周期内重建这些封装；基质可供性决定什么样的残差容易被吸收，皮肤、肌肉、肌腱、神经、痛觉和自修复组织，本身适合柔顺接触、连续反馈、损伤感知和局部修复，金属、电机、齿轮、硅片、线缆和刚性结构则更擅长精确、重复和高功率输出，却不天然具备生命材料的低功耗柔顺、自愈和连续感知。人用时间和生命基质隐藏复杂性，机器常常用工程系统展现复杂性。

这里所说的分摊可以被写成一个记账关系，

\begin{equation}
\bar{c}=\frac{C}{N}
\end{equation}

式中，$C$ 是形成一项能力所沉淀的总代价，包含演化试错、个体发育校准和社会环境改造的累计投入，必须先折算到同一口径（能量、时间或个体数）；$N$ 是这项能力在其使用史中被调用的总次数，由个体一生的行动次数与种群世代数的乘积给出，是一个无量纲计数；$\bar{c}$ 因此是单次调用分摊到的代价，单位随 $C$ 所选的口径而定。这个式子的含义在于，人类行动的轻巧不来自分子小，而来自分母大：$C$ 很大，$N$ 更大，因此 $\bar{c}$ 可以小到不进入当下的决策。机器要在短工程周期内重建同样的封装，$N$ 远小于演化与发育给出的分母，同样的 $C$ 摊下来就重得多。适用条件是一次性沉没与持续调用，只有当 $C$ 已经沉没且该项能力仍被反复调用时，摊销才成立；若环境改变使这项能力失效，原有的 $N$ 不再累积，摊销随之中断，$C$ 转为沉没成本。这个式子是记账方式，不给出 $C$ 与 $N$ 的实测值，也不表示生命身体在每一项能力上都更便宜。

但这并不意味着生命身体在所有维度上优越。机器在高精度重复、高载荷、高速加工、高温环境、强辐射场景和标准化制造中，常常远超生物身体，问题不在于生命材料天然“更高级”，而在于开放世界的柔顺接触、连续反馈、损伤感知、局部恢复和长时程行动，恰好是生命基质长期适配过的能力组合，机器可以在许多单项能力上超过身体，却仍然需要为这种组合能力重新设计封装结构。人类身体的轻巧还来自长期校准出的自状态先验：正常情况下，人不需要在每个动作前重新证明自己的手还在何处、关节是否存在、肌肉是否响应、身体是否仍保持基本一致性，身体通过本体感觉、前庭系统、运动经验和内稳态机制，提供了一个持续更新但并不每次归零的自我状态置信区间；机器则必须把这种自状态可信度工程化为传感器标定、冗余测量、状态估计、异常检测和安全断言，它不是没有“自我状态”，而是这种状态必须持续被测量、同步、校验和记录，身体中的默认信任，在机器中表现为一整套验证结构。

这两条轴合起来，把物理智能重新接回技术黏滞的主轴。比特层可以快速生成模型、策略和仿真，原子层却必须承受接触、摩擦、惯性、材料和安全后果，生命层已经用极长时间生成了一种适合现实的身体；AI 试图从比特层直接进入原子层，并复制生命层长期摊销出的身体能力，就会遭遇速率上的根本不匹配，模型更新快，身体改造慢，环境共适配更慢。物理 AI 的困难，深层上不是机器还不够聪明，而是快工程试图承接慢演化的历史积累，这一条也解释了为什么现实中的工程出路总是分成几种，而不能收束成一种。第一条是改造环境，减少残差来源，封闭园区、仓储、工厂、港口、专用道路、标准化接口，本质上都在把现实变得更可感知、更可预测、更少接触歧义，不是让机器理解整个开放世界，而是让世界局部结构化。第二条是改造身体，增强残差吸收，柔顺材料、触觉皮肤、灵巧手、弹性关节、欠驱动结构、低延迟本体反馈，都在尝试让机器拥有更好的物理承接层，人形机器人的吸引力之一，在于大量既有环境、工具和交互接口已经围绕人类身体尺度形成，但它也迫使机器承担双足平衡、全身协调、灵巧操作、安全接触和长期身体经验的高复杂度。第三条是增强模型，推远能力边界，更好的大脑可以用较少输入推断更多状态，可以提前预测错误，可以让现有硬件做更多事，但这不是逃逸，而是转移，复杂性从身体、传感器或环境转移到模型、数据、训练、仿真和验证中。第四条是收缩任务，并用运行体系兜底，很多现实可用的机器人系统，不是因为它们拥有通用身体，而是因为任务边界被精确限定，异常由人类或运行系统处理，机器处理被稳定化的常态，人类或运行体系处理超出边界的残差，这不是失败，但说明自主性并没有免费获得，只是被系统重新分配。

由此看，物理 AI 的成熟，不是单纯模型能力增长，也不是单纯传感器堆叠，而是建立足够可靠的\term{残差管理体系}，这个体系可以由身体、模型、传感器、控制器、环境设计、任务边界和运行体系共同组成，不同方案不是简单优劣，而是残差在不同位置被吸收，复杂性在不同位置展现。更进一步说，成熟的物理智能未必追求与开放现实发生无限高维接触，恰恰相反，它往往需要主动限制接触维度，划定可承接的物理边界，结构化环境、标准化工具、受限任务、柔顺身体、安全控制、远程接管和运行审计，都是在为机器划出一个可行动的容错区域；在这个区域内，残差可以存在，延迟可以存在，误差可以存在，但它们不会轻易穿透身体、模型、控制和运行体系的多重吸收结构。真正可靠的物理智能，不是让机器毫无边界地理解整个世界，而是让机器与世界共同形成一个可恢复、可验证、可维护的行动生态位。

形态计算之后，真正的问题不是身体是否参与计算，而是身体如何把残差跨时间尺度地封装起来，以及机器如何在缺少这种历史摊销与生命基质时重新承担残差。人不是因为拥有完整世界模型才行动，而是因为身体允许模型不完整；机器不是不能学习人的行动，也不是必须无限堆叠传感器，而是必须为缺失的身体封装寻找残差修正机制。物理智能的第一性难题就在这里：真正稀缺的不是抽象推理本身，而是一个经过漫长时间摊销、并且由合适基质承载的残差吸收结构。能不能进入开放世界，取决于系统是否拥有与它的介质相匹配的残差吸收结构。但这个结论并没有把问题说完，因为它把太多不同性质的工作都装进了“残差吸收”这一个说法里：展开可能性的工作可以在符号空间里低成本重试，封装风险的工作必须把自由度压进安全边界，显影错误的工作则要以证据链和长期观测为前提，三者需要的资源结构、时间尺度与确定性水平并不相同。如果残差的吸收位置决定了系统的重量，那么落在不同位置上的残差，是否本来就要求智能分化成不同的模态？
